\chapter{The Idea and History of Limits}

In the preceding chapters, we assembled the elementary cast of functions --
polynomials, exponentials, logarithms, and trigonometric curves -- by evaluating
them at concrete points. But what happens when the exact instant we care about
cannot be evaluated, collapsing into the forbidden ratio $0/0$?

There is a single idea sitting underneath every formula in this book that
resolves this riddle. It is not the derivative, and it is not the integral. It is
older than both, and both are defined in terms of it. It is the \emph{limit}.

A limit answers a question that sounds almost paradoxical: \emph{what value is a quantity heading
toward, even if it never actually gets there?} How
can we talk about where something is going if it never arrives? And yet your intuition handles
this kind of reasoning every day. When you watch a car decelerate toward a stop sign, you can
tell it is heading for a complete stop long before its wheels stop turning. When you pour
water into a glass and the level rises more and more slowly, you can see it approaching the
brim. The limit is the mathematics of \emph{approach} -- of destinations inferred from the
journey.

This chapter is about the \emph{idea} of a limit and the long history of how
mathematicians fought to understand it. The famous ``epsilon--delta'' definition can wait for
the next chapter. Here we want something
harder to get from a formula: we want you to \emph{see} what a limit is,
why calculus cannot live without it, and why it took roughly two hundred years and some of the
sharpest minds in history to say precisely what it means.

\section{A quantity that approaches}

Start with the most concrete possible picture. Consider the function
\[
  f(x) = \frac{x^2 - 1}{x - 1}.
\]
There is something wrong with this function at exactly one point. If you try to evaluate it at
$x = 1$, you get
\[
  f(1) = \frac{1 - 1}{1 - 1} = \frac{0}{0},
\]
which is undefined. Division by zero is not allowed, and $0/0$ is not allowed twice over. So
the function $f$ has a \emph{hole} at $x = 1$: it is defined everywhere else on the number
line, but at $x = 1$ it simply has no value.

Now here is the question that started calculus. The function is undefined \emph{at} $x = 1$,
but what is it doing \emph{near} $x = 1$? What value is it approaching as $x$ creeps toward
$1$?

\begin{labbox}
Before reading on, build the table yourself with a calculator. Compute $f(x) = \dfrac{x^2-1}{x-1}$
for $x = 0.9,\ 0.99,\ 0.999$ and then for $x = 1.1,\ 1.01,\ 1.001$. Watch what happens to the
output. Do not simplify the formula first -- plug the numbers in raw, the way a machine would.
\end{labbox}

Here is the table you should have produced. The values of $x$ march in toward $1$ from below
(the left) and from above (the right), and we watch the outputs.

\begin{center}
\begin{tabular}{rr@{\qquad}rr}
\toprule
$x$ (from left) & $f(x)$ & $x$ (from right) & $f(x)$ \\
\midrule
$0.9$    & $1.9$    & $1.1$    & $2.1$ \\
$0.99$   & $1.99$   & $1.01$   & $2.01$ \\
$0.999$  & $1.999$  & $1.001$  & $2.001$ \\
$0.9999$ & $1.9999$ & $1.0001$ & $2.0001$ \\
\bottomrule
\end{tabular}
\end{center}

The pattern is unmistakable. As $x$ approaches $1$ from either side, $f(x)$ approaches $2$. It
never \emph{equals} $2$ in this table, and it never can equal anything at $x=1$, because $f$ is
not defined there. But the outputs are clearly homing in on $2$. We say:
\[
  \lim_{x \to 1} \frac{x^2 - 1}{x - 1} = 2.
\]
Read this aloud as ``the limit, as $x$ approaches $1$, of $(x^2-1)/(x-1)$, equals $2$.''

\begin{focusbox}
The notation $\displaystyle\lim_{x\to a} f(x) = L$ is a statement about the \emph{approach},
not about the destination point itself. It says: \emph{as $x$ gets close to $a$ (but not equal
to $a$), the outputs $f(x)$ get close to $L$.} Whether $f$ is defined at $a$, and what value it
might have there, is a separate question that the limit deliberately ignores.
\end{focusbox}

Why does $f$ approach $2$? Algebra reveals the reason. As long as $x \neq 1$, we may factor and
cancel:
\[
  f(x) = \frac{x^2 - 1}{x - 1} = \frac{(x-1)(x+1)}{x - 1} = x + 1.
\]
The cancellation is legal precisely \emph{because} $x \neq 1$, so the factor $x - 1$ in the
denominator is not zero and we are allowed to divide it out. For every $x$ except $1$,
$f(x)$ simplifies directly to $x + 1$. The linear function $x + 1$ evaluates to $1 + 1 = 2$ at $x = 1$.
Consequently, as $x$ approaches $1$, $f(x)$ approaches $2$.

\begin{figure}[htbp]
\centering
\begin{tikzpicture}
\begin{axis}[
  width=11cm, height=7cm,
  title style={font=\small},
  xlabel={$x$}, ylabel={$f(x)$},
  xmin=-0.5, xmax=2.5, ymin=0, ymax=3.6,
  axis lines=left,
  samples=2, domain=-0.5:2.5,
  legend pos=north west, legend cell align=left,
  legend style={inner xsep=4pt},
  enlargelimits=false,
]
\addplot[curveblue, very thick, domain=-0.5:2.5] {x+1};
\addlegendentry{$f(x)=x+1$ (for $x\neq 1$)}
\addplot[mark=*, mark options={fill=white, draw=tangentred, line width=1.2pt}, only marks]
  coordinates {(1,2)};
\addplot[gray, dashed] coordinates {(1,0) (1,2)};
\addplot[gray, dashed] coordinates {(-0.5,2) (1,2)};
\node[tangentred, font=\footnotesize, right] at (axis cs:1.05, 2.0) {missing point $(1,2)$};
\end{axis}
\end{tikzpicture}
\caption{The graph of $f(x)=\dfrac{x^2-1}{x-1}$ is the straight line $y=x+1$ with a single
point punched out at $(1,2)$. The hollow circle marks the missing point. As $x$ slides toward
$1$ along the line, the height slides toward $2$ -- regardless of the fact that the point
itself is absent.}
\label{ch05:hole}
\end{figure}

\begin{intuitionbox}
A limit is what your eye does automatically when it follows a curve toward a point. If you trace
the line in Figure~\ref{ch05:hole} with your finger heading for $x=1$, your finger arrives at
height $2$ -- and the tiny hole does not stop you, because your finger is reading the
\emph{trend} of the curve, not the single missing pixel. The limit is the trend.
\end{intuitionbox}

\begin{recoverybox}
\textbf{Pause and think:} Suppose someone argues that $\lim_{x\to 1} \dfrac{x^2-1}{x-1}$ cannot equal $2$
because ``when you plug in $x=1$, you get $0/0$, which is undefined.'' How would you answer in a single
sentence? The limit does not ask what happens \emph{at} the point; it asks what value the outputs are
marching toward as you get arbitrarily close to the point.
\end{recoverybox}

\subsection{Approaching the undefined point}

You might shrug and say: this is a fuss over nothing. Just use the line $x+1$ and forget the
hole. Yet that hole shows why limits are needed in the first place.
Consider what happens throughout calculus.

When we measure the steepness of a curve at a point, we will compute the slope of a line
through two points on the curve and then let those two points slide together. At the instant
they coincide, the slope formula becomes
\[
  \frac{\text{rise}}{\text{run}} = \frac{0}{0},
\]
because the two points are the same point: zero rise over zero run. When we measure the area
under a curve, we will add up the areas of thin rectangles and let their width shrink to zero;
in the limit each rectangle has zero width, and we are summing infinitely many zeros. Every
central computation in calculus runs straight into a $0/0$ or an ``infinity of zeros,'' an
expression that is meaningless if you plug in the forbidden value directly. The limit is the
tool that lets us extract a sensible, finite answer from these forbidden expressions by asking
not ``what is the value \emph{at} the bad point?'' but ``what is the value \emph{approaching} as
we close in on the bad point?''

\begin{aestheticbox}
This is the deep reason the limit is the foundation. Calculus is the study of quantities at an
\emph{instant} -- instantaneous speed, instantaneous slope, the area swept out at a precise
moment. But an instant has no duration, and at an instant the naive formulas collapse to $0/0$.
The limit is how we make sense of ``at an instant'' without ever actually dividing by zero. It
is the bridge from the smooth, continuous world we want to describe to the arithmetic of finite
numbers we are allowed to do.
\end{aestheticbox}

\section{A limit that does not need a hole}

It would be a mistake to think limits are only about repairing broken functions. Most of the
time the limit of a function at a point is simply its value there, and that is the boring,
reassuring case. If $g(x) = x^2 + 3$, then
\[
  \lim_{x \to 2} g(x) = 2^2 + 3 = 7,
\]
and the outputs near $x = 2$ approach $7$ for the unremarkable reason that $g$ is a nice smooth
function with $g(2) = 7$. There is no hole, no drama. The function approaches the value it
actually has (Figure~\ref{ch05:nohole}).

\begin{figure}[htbp]
\centering
\begin{tikzpicture}
\begin{axis}[
  width=10cm, height=6cm,
  title style={font=\small},
  title={Continuity: limit equals function value},
  xlabel={$x$}, ylabel={$g(x)$},
  xmin=0, xmax=3.2, ymin=2.5, ymax=11,
  axis lines=left,
  samples=80, domain=0:3.2,
  enlargelimits=false,
]
\addplot[curveblue, very thick, domain=0:3.2] {x^2+3};
\addplot[mark=*, mark options={fill=curveblue, draw=curveblue}, only marks]
  coordinates {(2,7)};
\addplot[gray, dashed] coordinates {(2,2.5) (2,7)};
\addplot[gray, dashed] coordinates {(0,7) (2,7)};
\node[curveblue, font=\footnotesize, anchor=south east] at (axis cs:1.9, 7.1) {point $(2,7)$};
\end{axis}
\end{tikzpicture}
\caption{The graph of $g(x)=x^2+3$ has no hole at $x=2$. As $x$ slides toward $2$ the height
slides toward $7$ -- and the point $(2,7)$ is genuinely \emph{there} (filled dot), so here the
limit and the value agree. Contrast this with the punched-out point of Figure~\ref{ch05:hole}.}
\label{ch05:nohole}
\end{figure}

Functions that behave this well -- where the limit at a point equals the value at the point --
are called \emph{continuous} (informally, curves you can sketch without lifting your pencil from
the page), and they get their own chapter. For now, just notice that the
limit and the value can either \emph{agree} (the continuous case, $g$ above) or \emph{disagree}
(the hole case, $f$ above, where the limit is $2$ but the value does not exist). The limit
concept earns its place precisely because it does not care which case you are in. It always asks the
same question -- where are the outputs heading? -- and that question has an answer of $2$ for
$f$ at $x=1$ and an answer of $7$ for $g$ at $x=2$, whether or not the function bothers to be
defined and equal there.

\begin{happeningbox}
When you write $\lim_{x\to a} f(x) = L$, you are making a claim about \emph{infinitely many}
values of $f$ at once -- all the values at points near $a$. The claim is: \emph{you can force
the outputs to be as close to $L$ as you like, just by keeping $x$ close enough to $a$.} The
phrase ``as close as you like'' is the heart of it. It is not enough that the outputs get
\emph{somewhat} close; they must be squeezable arbitrarily close. We deliberately exclude
$x = a$ itself from the conversation, which is why the value $f(a)$ -- present, absent, or
weird -- has no vote. The next chapter turns ``as close as you like'' into the precise
$\eps$--$\delta$ machinery. Everything there is just this sentence made unambiguous.
\end{happeningbox}

\section{When left and right disagree: one-sided limits}

So far our examples approached the same value whether $x$ came in from the left or from the
right. That symmetry is not guaranteed, and watching it break is one of the most instructive
things you can do.

Consider the \emph{sign-step} function
\[
  H(x) = \begin{cases} 0 & x < 0, \\[2pt] 1 & x \ge 0. \end{cases}
\]
This is the Heaviside step (named after electrical engineer Oliver Heaviside, it models an abrupt
switch flipped from ``off'' to ``on''). It is flat at height $0$ for all negative $x$, then jumps
suddenly to height $1$ at $x = 0$ and stays there. Ask the limit question at $x = 0$: what is $H$
approaching?

\begin{figure}[htbp]
\centering
\begin{tikzpicture}
\begin{axis}[
  width=11cm, height=5.5cm,
  title style={font=\small},
  title={Jump: one-sided limits disagree},
  xlabel={$x$}, ylabel={$H(x)$},
  xmin=-2.2, xmax=2.2, ymin=-0.4, ymax=1.6,
  axis lines=left, ytick={0,1},
  enlargelimits=false,
]
\addplot[curveblue, very thick, samples=2, domain=-2.2:0] {0};
\addplot[curveblue, very thick, samples=2, domain=0:2.2] {1};
\addplot[mark=*, mark options={fill=white, draw=curveblue, line width=1pt}, only marks]
  coordinates {(0,0)};
\addplot[mark=*, mark options={fill=curveblue, draw=curveblue}, only marks]
  coordinates {(0,1)};
\addplot[gray, dashed] coordinates {(0,-0.4) (0,1.6)};
\node[curveblue, font=\footnotesize, anchor=south] at (axis cs:-1.1, 0.08) {left approach: $\to 0$};
\node[curveblue, font=\footnotesize, anchor=south] at (axis cs:1.1, 1.08) {right approach: $\to 1$};
\end{axis}
\end{tikzpicture}
\caption{The Heaviside step $H(x)$. Approaching $x=0$ from the left, the outputs sit at $0$;
approaching from the right, they sit at $1$. The two approaches disagree, so there is no single
value that $H$ is heading toward.}
\label{ch05:heaviside}
\end{figure}

Figure~\ref{ch05:heaviside} shows the jump. If you walk in toward $0$ from the negative side --
$x = -0.1, -0.01, -0.001$ -- every output is exactly $0$. The function is approaching $0$. But if you walk in from the positive side --
$x = 0.1, 0.01, 0.001$ -- every output is exactly $1$. The function is approaching $1$. The two
sides give different answers. There is no single number that $H(x)$ is ``heading toward'' as $x
\to 0$, because it depends entirely on which direction you came from.

We capture this with \emph{one-sided limits}. The \emph{left-hand limit} watches only $x$
coming from below:
\[
  \lim_{x \to 0^-} H(x) = 0,
\]
where the little minus sign on $0^-$ means ``$x$ approaches $0$ through values less than $0$.''
The \emph{right-hand limit} watches only $x$ coming from above:
\[
  \lim_{x \to 0^+} H(x) = 1,
\]
with $0^+$ meaning ``through values greater than $0$.'' Both one-sided limits exist and are
perfectly definite. They simply disagree.

\begin{focusbox}
The ordinary (two-sided) limit $\lim_{x\to a} f(x)$ exists \emph{only when both one-sided limits
exist and are equal}. In symbols,
\[
  \lim_{x \to a} f(x) = L
  \quad\Longleftrightarrow\quad
  \lim_{x \to a^-} f(x) = L \ \text{ and }\ \lim_{x \to a^+} f(x) = L.
\]
For the Heaviside step, the left limit is $0$ and the right limit is $1$; they disagree, so the
two-sided limit $\lim_{x\to 0} H(x)$ \emph{does not exist}. ``Does not exist'' is a legitimate,
useful answer, not a failure.
\end{focusbox}

\begin{intuitionbox}
Think of two hikers approaching a canyon from opposite rims. The left hiker ends at the canyon
floor; the right hiker ends on the far rim. If someone asks ``where does the trail arrive?''
there is no single answer -- it depends which hiker you follow. A two-sided limit exists only
when both hikers, coming from opposite directions, arrive at the \emph{same} place. One-sided
limits let us still describe each hiker's destination separately.
\end{intuitionbox}

\begin{recoverybox}
\textbf{Pause and think:} If you are told that $\lim_{x\to 3^-} f(x) = 5$ and $\lim_{x\to 3^+} f(x) = 5$,
can you conclude that $f(3) = 5$? No! The two-sided limit exists and equals $5$, but $f(3)$ could be
undefined (a hole) or equal to $100$ (a displaced point). The limit describes the approach; the function
value is whatever sits at the destination itself.
\end{recoverybox}

\begin{applicationbox}
One-sided limits are the natural way to model \emph{sudden
events}. Flip a switch and the voltage jumps; the current ``before'' and ``after'' have
different limits at the instant of the flip. In economics, a tax bracket creates a jump in the
marginal rate: the rate just below the threshold and just above it are different one-sided
limits. In probability, the cumulative distribution of a discrete variable steps up at each
value, exactly like the Heaviside function. Whenever reality has a genuine jump, the two-sided
limit fails and the one-sided limits carry the real information.
\end{applicationbox}

\subsection{A more violent disagreement}

A jump is a polite failure: each side still settles on a number. Limits can fail more
dramatically. Consider $p(x) = 1/x$ near $x = 0$. From the right, as $x = 0.1, 0.01, 0.001$, the
outputs are $10, 100, 1000$ -- shooting off to $+\infty$. From the left, $x = -0.1, -0.01,
-0.001$ give $-10, -100, -1000$ -- plunging to $-\infty$. Neither one-sided limit settles on a
finite number; the function values grow without bound, as shown in Figure~\ref{ch05:reciprocal}. We write
$\lim_{x\to 0^+} 1/x = +\infty$ and $\lim_{x\to 0^-} 1/x = -\infty$ as shorthand for ``grows
without bound,'' but you should hear those as descriptions of \emph{how} the limit fails, not as
ordinary numerical answers. We will treat such infinite behavior carefully in the chapter on
asymptotes.

\begin{figure}[htbp]
\centering
\begin{tikzpicture}
\begin{axis}[
  width=10cm, height=6.5cm,
  title style={font=\small},
  title={Infinite limits: blow-up as $x \to 0$},
  xlabel={$x$}, ylabel={$1/x$},
  xlabel style={anchor=south west},
  ylabel style={anchor=south east},
  xmin=-2.4, xmax=2.4, ymin=-6, ymax=6,
  xtick={1,2},
  extra x ticks={-2,-1},
  extra x tick labels={},
  ytick={-4,-2,2,4},
  axis lines=middle,
  samples=200, restrict y to domain=-6:6,
  enlargelimits=false,
]
\node[font=\footnotesize, above] at (axis cs:-1, 0) {$-1$};
\node[font=\footnotesize, above] at (axis cs:-2, 0) {$-2$};
\addplot[curveblue, very thick, domain=0.17:2.15] {1/x};
\addplot[curveblue, very thick, domain=-2.15:-0.17] {1/x};
\end{axis}
\end{tikzpicture}
\caption{$p(x)=1/x$. The right-hand branch races up toward $+\infty$ as $x\to 0^+$; the
left-hand branch dives toward $-\infty$ as $x\to 0^-$. There is no finite value being
approached from either side.}
\label{ch05:reciprocal}
\end{figure}

\section{The historical crisis: a triumph built on sand}

Now for the history, where the idea of a limit turns dramatic. Calculus was invented --
twice, independently -- in the late seventeenth century. Isaac Newton in England (working in
the 1660s and 1670s, publishing later) and Gottfried Wilhelm Leibniz in Germany (publishing in
the 1680s) each built a calculus of \emph{infinitesimals} (idealized quantities so tiny that
they are smaller than any positive fraction, yet still greater than zero). Their methods were
powerful. Newton used his ``fluxions'' (his term for the instantaneous rates of
change of flowing quantities, from the Latin for flow) to derive the laws of planetary motion;
Leibniz's elegant $\dd y/\dd x$ and $\int$ notation (an elongated ``S'' for sum) -- which we
still use -- let problems that had stumped the ancients fall in a few lines. By any practical
measure, calculus was a colossal success.

And yet at its logical core sat something nobody could justify.

To see the puzzle, picture how both Newton and Leibniz computed the slope of a curve.
Consider $y = x^2$ at an arbitrary point $x$. You pick a nearby point by letting $x$ change by
a tiny amount $o$ (Newton's notation) or $\dd x$ (Leibniz's), compute the resulting change in
$y$, and form the ratio of rise to run (Figure~\ref{ch05:fig-secant-tangent}).

\begin{figure}[htbp]
\centering
\begin{tikzpicture}
\begin{axis}[
  width=10.5cm, height=6.8cm,
  title style={font=\small},
  title={Secant line approaching tangent line on $y=x^2$},
  xlabel={$x$}, ylabel={$y$},
  xmin=0, xmax=2.3, ymin=0, ymax=4.8,
  axis lines=left,
  legend pos=north west,
  legend cell align=left,
  legend style={font=\footnotesize},
  enlargelimits=false,
]
\addplot[curveblue, very thick, domain=0:2.18, samples=60] {x^2};
\addlegendentry{$y=x^2$}
\addplot[tangentred, thick, dashed, domain=0.4:2.2] {2*x - 1};
\addlegendentry{tangent (slope $= 2$)}
\addplot[quotientorange, thick, domain=0.5:2.1] {2.6*x - 1.6};
\addlegendentry{secant (slope $= 2+o$)}
\addplot[mark=*, mark options={fill=curveblue, draw=curveblue}, only marks]
  coordinates {(1,1)};
\addplot[mark=*, mark options={fill=quotientorange, draw=quotientorange}, only marks]
  coordinates {(1.6,2.56)};
\addplot[gray, dashed] coordinates {(1,0) (1,1)};
\addplot[gray, dashed] coordinates {(1.6,0) (1.6,2.56)};
\node[curveblue, font=\footnotesize, anchor=east] at (axis cs:0.95, 1.25) {$P(1,1)$};
\node[quotientorange, font=\footnotesize, anchor=south east] at (axis cs:1.55, 2.7) {$Q(1+o, (1+o)^2)$};
\draw[-{Stealth[length=2.5mm]}, quotientorange, thick] (axis cs:1.5, 2.25) to[bend left=15] (axis cs:1.15, 1.35);
\node[quotientorange, font=\footnotesize, anchor=west] at (axis cs:1.35, 1.7) {$o \to 0$};
\end{axis}
\end{tikzpicture}
\caption{The geometric puzzle that troubled early calculus. To find the slope of $y=x^2$ at
$P(1,1)$, we pick a second point $Q(1+o, (1+o)^2)$ and connect them with an orange secant line
(slope $2+o$). As $o$ shrinks toward zero, $Q$ slides along the curve toward $P$, and the secant
line swings into the dashed red tangent line (slope $2$). For any nonzero $o$, the points are
distinct and the division is valid; the limit is the slope the secant is heading toward.}
\label{ch05:fig-secant-tangent}
\end{figure}

For $y = x^2$, the algebra runs:
\[
  \frac{(x+o)^2 - x^2}{o}
  = \frac{x^2 + 2xo + o^2 - x^2}{o}
  = \frac{2xo + o^2}{o}
  = 2x + o.
\]
Then came the suspicious maneuver: you \emph{discard} the leftover $o$ to conclude that the
instantaneous slope is $2x$.

Notice the problem in that calculation. In the third equality we \emph{divided} by $o$, which is
only legal if $o \neq 0$. Then in the final step we \emph{threw away} the term $o$, which is
only legal if $o = 0$. The increment $o$ had to be nonzero to allow division, and zero to
disappear from the answer. Newton and Leibniz both sensed this tension; neither could explain
in clean logical terms why the answer $2x$ was nonetheless indisputably correct.

This was not a minor bookkeeping quirk. It was a contradiction at the very base of the most
celebrated mathematics of the Enlightenment, and it did not escape notice. The most famous
attack came not from a rival mathematician, but from a philosopher and clergyman: Bishop
George Berkeley.

\begin{historybox}
In 1734, the philosopher Bishop George Berkeley published \emph{The Analyst},
criticizing the logical foundations of early calculus. He pointed out that in calculating
derivatives, the increment $o$ was treated as nonzero in order to divide by it, and then
treated as zero to eliminate it from the result. In a famous passage, he challenged mathematicians:
\begin{quote}
\emph{``And what are these fluxions? The velocities of evanescent increments? And what are these same evanescent increments? They are neither finite quantities, nor quantities infinitely small, nor yet nothing. May we not call them the ghosts of departed quantities?''}
\end{quote}
Berkeley demonstrated that while the results were correct, the methods lacked rigorous justification.
This critique motivated nineteenth-century mathematicians to establish limits on formal foundations
without relying on infinitesimals. In the 1960s, Abraham Robinson's nonstandard analysis proved
that infinitesimals can also be formulated with complete mathematical rigor.
\end{historybox}

\begin{walkbox}
Take Berkeley's challenge seriously on your next walk. The computation $\frac{(x+o)^2-x^2}{o} =
2x + o$ is exact and innocent as long as $o\neq 0$. The whole crisis lives in a single move: at
the end we replace $2x + o$ by $2x$. Ask yourself: \emph{is that an approximation, or is it
exact?} If it is an approximation, calculus is only approximately true -- but it is not, the
answers are exact. If it is exact, then in what sense is dropping a nonzero $o$ legitimate? Hold
that question. The limit is precisely the answer, and feeling the tension first will make the
resolution land.
\end{walkbox}

\subsection{What the early masters got right -- and what was missing}

Do not read this history as ``Newton and Leibniz were sloppy.'' They were
not sloppy; they were \emph{early}. Their physical and geometric intuition was nearly flawless.
When Newton spoke of a quantity's ``ultimate ratio,'' he described the ratio $2x+o$ \emph{as $o$
vanishes} -- not its value at any fixed nonzero $o$, and not a meaningless ratio at $o=0$, but
the value the ratio is \emph{heading toward}. Read that sentence again: Newton was reaching, in
words, for exactly the limit concept. He wrote of ``the limit to which they always approach, and
to which they approach nearer than by any given difference, but never go beyond.'' That is almost
a definition of a limit -- in 1687.

The gap was language, not insight. What Newton and Leibniz lacked was not the geometric picture, but a way to state what ``approaching'' meant without invoking infinitely small numbers that were somehow both nonzero and zero. Once nineteenth-century mathematicians defined the limit cleanly, every formula Newton and Leibniz had found stayed intact -- only the foundation underneath changed.

\begin{happeningbox}
What is really happening in the $0/0$ derivative computation? \emph{Nothing is ever divided by
zero, and nothing nonzero is ever discarded.} The expression $\frac{(x+o)^2-x^2}{o}$ equals
$2x+o$ for \emph{every} nonzero $o$ -- this is exact algebra, no approximation. The slope is then
\emph{defined} as the limit of that exact expression as $o\to 0$, i.e.\ the number that $2x+o$ is
heading toward, which is $2x$. We never set $o=0$; we ask where $2x+o$ goes as $o$ shrinks.
No infinitesimal number is needed: ``$o\to 0$'' represents a limiting process, not an arithmetic
quantity.
\end{happeningbox}

\section{Cauchy and Weierstrass: rebuilding on rock}

For over a century after Newton's \emph{Principia}, calculus advanced
spectacularly while its logical foundation remained incomplete. Leonhard Euler
and his contemporaries computed extensively with infinite series and infinitesimals,
guided by strong physical and geometric intuition. But as functions got
stranger and infinite series got subtler, intuitive methods occasionally produced
contradictions, and establishing clear definitions became essential to determine when
manipulations were mathematically sound.

In the 1820s, Augustin-Louis Cauchy, lecturing at the \'Ecole Polytechnique in Paris, made a
decisive choice in his \emph{Cours d'analyse} (1821): he made the \emph{limit} the primary
concept and defined everything else -- continuity, the derivative, the integral, and convergence
(the property of an infinite sequence settling on a single destination) -- in terms of it. He
wrote that a variable's limit is the fixed value it approaches so that it ends up differing from
it by ``as little as one wishes.''

In this formulation, an infinitesimal is no longer treated as a fixed nonzero quantity
smaller than any real number, but as a variable that converges toward zero.

Cauchy's phrase ``as little as one wishes'' is the seed of everything. It contains a hidden
challenge-and-response: \emph{however} small a gap you want, I can guarantee the difference is
smaller, provided the variable is far enough along. But Cauchy still leaned on intuitive verbal
phrasing -- ``approaches indefinitely,'' ``successive values'' -- and his words occasionally hid
subtle errors (for instance, his famous assertion that any convergent sum of continuous
functions is continuous is, without extra conditions, false).

The final tightening fell to Berlin. In the 1850s through 1870s, Karl Weierstrass, a former
schoolteacher who became the leading figure in mathematical analysis (the branch of mathematics
devoted to limits and continuous processes), removed the last trace of motion and mystery from
the limit. Where Cauchy still spoke of a variable \emph{approaching} -- a word that suggests time
and movement -- Weierstrass replaced the whole idea with a static inequality between two ordinary positive numbers, $\eps$ (output tolerance) and $\delta$ (input tolerance): for every $\eps > 0$, there exists $\delta > 0$ such that $0 < |x-a| < \delta \implies |f(x)-L| < \eps$. No motion, no infinitesimals -- just distances on the number line.

\begin{aestheticbox}
Notice the direction of the centuries-long motion: from \emph{dynamic} to \emph{static}. Newton's
quantities \emph{flowed}; Cauchy's variables \emph{approached}; Weierstrass's $\eps$ and $\delta$
just \emph{sit there} in a logical relationship with no motion at all. This is a recurring pattern
in mathematics -- we tame a slippery dynamic intuition by reformulating it as a static statement
about ``for all $\dots$ there exists $\dots$.'' Trading the feeling of motion for a static inequality can feel austere at first, but it gives us a definition sharp enough to prove theorems without guesswork.
\end{aestheticbox}

\begin{figure}[htbp]
\centering
\begin{tikzpicture}[
  node distance=0.8cm and 1.0cm,
  card/.style={draw=#1, fill=#1!7, rounded corners=3pt, line width=0.8pt, inner sep=6pt, text width=5.6cm, align=left, font=\footnotesize},
  heading/.style={font=\bfseries\footnotesize, text=#1},
  conn/.style={-Stealth, thick, gray!70, shorten <=2pt, shorten >=2pt}
]
\node[card=curveblue] (newton) {
  \textbf{\color{curveblue} 1. Newton \& Leibniz (1670s--1680s)}\\[2pt]
  \emph{Dynamic Infinitesimals}\\
  Variables flow continuously; increments $\dd x$ and $o$ are nonzero during division, then dropped as zero at the end:
  $\frac{2xo + o^2}{o} = 2x+o \longrightarrow 2x$.
};

\node[card=tangentred, below=0.7cm of newton] (berkeley) {
  \textbf{\color{tangentred} 2. Berkeley's Critique (1734)}\\[2pt]
  \emph{Critique of Infinitesimals}\\
  Attacked the logical contradiction: $o$ cannot be nonzero to divide and zero to discard in the same calculation without formal justification.
};

\node[card=quotientorange, right=0.8cm of newton] (cauchy) {
  \textbf{\color{quotientorange} 3. Cauchy (1820s)}\\[2pt]
  \emph{Dynamic Limits}\\
  Made the limit primary. The infinitesimal is demoted from a mysterious number to a \emph{variable tending to zero}; difference can be made ``as small as one wishes.''
};

\node[card=structurepurple, below=0.7cm of cauchy] (weierstrass) {
  \textbf{\color{structurepurple} 4. Weierstrass (1860s)}\\[2pt]
  \emph{Static Quantifiers: $\eps$--$\delta$}\\
  Motion and time eliminated completely. Rigorous static guarantee: for every tolerance $\eps > 0$, there exists $\delta > 0$ such that $0 < |x-a| < \delta \implies |f(x)-L| < \eps$.
};

\draw[conn] (newton) -- (berkeley) node[midway, right, font=\tiny, text=gray!80] {challenged by};
\draw[conn] (berkeley) -- (cauchy) node[midway, above, sloped, font=\tiny, text=gray!80] {rebuilt via};
\draw[conn] (cauchy) -- (weierstrass) node[midway, right, font=\tiny, text=gray!80] {formalized as};
\end{tikzpicture}
\caption{The two-century conceptual progression of the limit: from physical intuition of flowing infinitesimals to a completely static arithmetic statement.}
\label{ch05:fig-timeline}
\end{figure}

\begin{recoverybox}
\textbf{Pause and think:} In everyday language, a ``limit'' is a boundary you cannot cross, like a
speed limit on a highway. In calculus, can a function cross its limit? Yes, repeatedly! Consider an
oscillating spring settling down, or $f(x) = \frac{\sin x}{x}$ as $x\to\infty$. The graph crosses the
limiting height of $0$ infinitely many times as it squeezes closer and closer. In calculus, a limit is a
destination being approached, not an impassable wall.
\end{recoverybox}

\subsection{Why intuition alone is insufficient}

You might reasonably ask: if the intuitive idea -- ``outputs head toward $L$'' -- gets the right
answers, why bother with the austere $\eps$--$\delta$ machinery at all? Why did this need fixing?

The answer is that intuition, unaided, eventually lies to you. Consider how you would
\emph{prove} a basic fact such as ``the limit of a sum is the sum of the limits,'' or decide a
genuinely hard case where the function oscillates infinitely fast as it approaches the point.
Intuition has no traction on these; it can only gesture. Worse, the nineteenth century discovered
functions that violate every naive expectation -- a function continuous at every point yet with a
sharp corner at every point, so that ``follow the curve with your finger'' becomes meaningless.
On such monsters, the only way to know what is true is to have a definition precise enough to
\emph{compute} with, one that does not quietly assume the function is as tame as the
pictures in this chapter.

\begin{trapbox}
The most common misconception about limits, and the one the whole history is fighting, is the
belief that $\lim_{x\to a} f(x)$ means ``plug in $x=a$.'' Sometimes plugging in works (the
continuous case), and that lucky coincidence breeds the bad habit. But the limit is fundamentally
a statement about values \emph{near} $a$ with $a$ \emph{itself excluded}. Our very first example,
$\frac{x^2-1}{x-1}$, has a perfectly good limit of $2$ at a point where you \emph{cannot} plug in
at all. Train yourself to read $\lim_{x\to a}$ as ``as $x$ approaches $a$,'' never as ``at
$x=a$.''
\end{trapbox}

\begin{whybox}
Why does turning ``approaches'' into a static $\eps$--$\delta$ statement actually work -- what
makes it the definitive formalization? It captures the
exact operational content of ``as close as you like.'' The phrase already contains a two-player
challenge: a skeptic challenges you with a tolerance (``get within $\eps = 0.001$ of $L$''), and you
must respond with a guarantee (``keep $x$ within $\delta$ of $a$ and I promise you will''). If you
can meet this challenge for \emph{every} $\eps$ -- no matter how small -- then in
every practical and provable sense the outputs are being forced toward $L$, and there is no
wiggle room left for them to be heading anywhere else. The genius of Weierstrass's formulation is
that it converts the vague ``gets closer and closer'' into a finite, checkable promise about two
ordinary positive numbers. The next chapter is entirely devoted to playing this game until it is
second nature.
\end{whybox}

\section{A numerical caution: tables suggest, they do not prove}

The table that opened this chapter was persuasive, and tables are a wonderful way to \emph{guess}
a limit. But a table can mislead, and seeing how is excellent training for why we eventually need
proof rather than evidence.

\begin{labbox}
Tabulate $f(x) = \sin(\pi/x)$ at the points $x = 1,\ \tfrac{1}{2},\ \tfrac{1}{3},\
\tfrac{1}{4},\ \dots$, i.e.\ $x = 1/n$ for $n=1,2,3,\dots$. At every one of these points $\pi/x =
n\pi$, so $\sin(\pi/x) = \sin(n\pi) = 0$. The table reads $0,0,0,0,\dots$ and screams ``the limit
as $x\to 0^+$ is $0$.'' Now tabulate instead at $x = 2,\ \tfrac{2}{5},\ \tfrac{2}{9},\ \dots$,
i.e.\ $x = 2/(4k+1)$. There $\pi/x = (4k+1)\pi/2$, an odd multiple of $\pi/2$, so the output is
$+1$ every time. A second table, just as correct, screams ``the limit is $1$.''
\end{labbox}

Both tables are arithmetically correct, and they contradict each other. The truth is that
$\sin(\pi/x)$ oscillates between $-1$ and $+1$ faster and faster as $x \to 0^+$, never settling
on anything: the limit \emph{does not exist}, as the crowding oscillations of
Figure~\ref{ch05:oscillation} show. No finite table could ever reveal this for certain,
because a table only samples finitely many points and the misbehavior hides in between them.

\begin{figure}[htbp]
\centering
\begin{tikzpicture}
\begin{axis}[
  width=12cm, height=5.5cm,
  title style={font=\small},
  title={Infinite oscillation near the origin},
  xlabel={$x$}, ylabel={$\sin(\pi/x)$},
  xmin=0, xmax=1.05, ymin=-1.4, ymax=1.4,
  axis lines=left, samples=200, domain=0.04:1.05,
  enlargelimits=false,
]
\addplot[curveblue, thick] {sin(deg(pi/x))};
\end{axis}
\end{tikzpicture}
\caption{$\sin(\pi/x)$ near $x=0$. The oscillations crowd together without bound as $x\to 0^+$,
so the curve never homes in on a single height. A numerical table that happens to land on the
zeros (or on the peaks) would falsely suggest a limit. Evidence is not proof.}
\label{ch05:oscillation}
\end{figure}

\begin{focusbox}
The moral is not ``never use tables.'' Tables and graphs are how you \emph{form} the conjecture,
and forming the right conjecture is most of the battle. The moral is that a finite amount of
evidence cannot \emph{settle} a statement about infinitely many points. That gap -- between
suggestive evidence and airtight proof -- is exactly the gap the $\eps$--$\delta$ definition was
built to close. This is why the next chapter exists.
\end{focusbox}

\section{Summary: The Core Concept of a Limit}

Here is the whole chapter in a single paragraph, the kind you could recite on a walk.

A limit is the value a function's outputs \emph{head toward} as the input is driven toward some
point, deliberately ignoring what -- if anything -- happens \emph{at} that point. It exists when
the approach from the left and the approach from the right agree on a single destination, and it
fails to exist when they disagree (a jump), or when the outputs run off to infinity, or when they
oscillate forever without settling. Newton and Leibniz used this idea to build calculus on top of ``infinitesimals'' whose logic Berkeley rightly challenged; Cauchy put the limit at the center of the subject as a variable shrinking toward zero; and Weierstrass finally stripped out all talk of motion, turning ``approaches'' into a precise challenge-and-response test with $\eps$ and $\delta$.

\begin{connectionbox}
Every major concept ahead is built on the limit. The \emph{derivative} is the limit of a slope as two
points merge ($0/0$ tamed). The \emph{integral} is the limit of a sum of shrinking rectangles
(infinitely many zeros tamed). \emph{Continuity} is the statement that a function's limit equals
its value. \emph{Infinite series} are limits of partial sums. \emph{Improper integrals} are
limits of ordinary ones. If you truly own the idea in this chapter -- a destination inferred from
an approach, with the endpoint itself excluded -- you own the spine of the entire subject, and
each later technique becomes a special way of computing one particular limit.
\end{connectionbox}

In Chapter~6, we turn Weierstrass's definition into a practical tool. We will make
$\eps$ and $\delta$ concrete by treating them as an input-output tolerance game.

\section{Exercises}

\begin{exercise}
(\emph{Mental check: limit versus value.}) Answer each question in one or two sentences:
\begin{enumerate}[label=(\alph*)]
  \item If you know that $\lim_{x\to 4} f(x) = 7$, can you conclude what $f(4)$ is? Could $f(4)$ be undefined? Could $f(4) = 99$?
  \item If you know that $f(4) = 7$, does that guarantee that $\lim_{x\to 4} f(x) = 7$? Sketch or describe a simple counterexample.
\end{enumerate}
\end{exercise}

\begin{exercise}
(\emph{Mental check: one-sided consistency.}) Suppose $f$ is a function defined near $x = 2$.
\begin{enumerate}[label=(\alph*)]
  \item If $\lim_{x\to 2^-} f(x) = 3$ and $\lim_{x\to 2^+} f(x) = 3$, must the two-sided limit $\lim_{x\to 2} f(x)$ exist? If so, what is its value?
  \item If instead $\lim_{x\to 2^-} f(x) = 3$ and $\lim_{x\to 2^+} f(x) = 3.0001$, does the two-sided limit exist? Why or why not?
\end{enumerate}
\end{exercise}

\begin{exercise}
Build a table of values for $\displaystyle g(x) = \frac{x^2 - 9}{x - 3}$ at $x = 2.9$, $2.99$,
$2.999$ and $x = 3.1$, $3.01$, $3.001$. Conjecture $\lim_{x\to 3} g(x)$ from the table, then confirm
your conjecture algebraically by factoring and cancelling. Explain in one sentence why the
cancellation is legitimate even though $g(3)$ is undefined.
\end{exercise}

\begin{exercise}
For each function, decide from a graph or short table whether the two-sided limit at the given
point exists, and if so what it is. If it does not exist, state both one-sided limits (or explain
that the behavior is infinite or oscillatory).
\begin{enumerate}[label=(\alph*)]
  \item $\displaystyle \frac{\abs{x}}{x}$ at $x = 0$.
  \item $\displaystyle \frac{x^3 - x}{x - 1}$ at $x = 1$.
  \item $\displaystyle \frac{1}{(x-2)^2}$ at $x = 2$.
\end{enumerate}
\end{exercise}

\begin{exercise}
Consider the piecewise function
\[
  f(x) = \begin{cases} x + 1 & x < 2, \\ 5 & x = 2, \\ 3x - 3 & x > 2. \end{cases}
\]
Compute $\lim_{x\to 2^-} f(x)$, $\lim_{x\to 2^+} f(x)$, and $\lim_{x\to 2} f(x)$. Then state the
value $f(2)$. Which of these four numbers can disagree with the others, and why does that not
contradict anything?
\end{exercise}

\begin{exercise}
(\emph{Berkeley's challenge.}) Carry out the early-calculus computation for
$y = x^3$. That is, expand $\dfrac{(x+o)^3 - x^3}{o}$, simplify (dividing by $o$ where $o\neq 0$),
and then identify the value the resulting expression heads toward as $o\to 0$. State clearly which
step requires $o\neq 0$ and which step requires $o\to 0$, and explain in your own words why no
single number $o$ is ever required to be both zero and nonzero.
\end{exercise}

\begin{exercise}
Explain, in plain prose and to a friend who has not taken calculus, why
$\displaystyle\lim_{x\to a} f(x)$ deliberately ignores the value $f(a)$. Give a specific example
where the limit exists but the function is undefined at the point, and a second example where the
limit exists, the function is defined at the point, but the two \emph{disagree}.
\end{exercise}

\begin{exercise}
(\emph{Tables can lie.}) Let $h(x) = \cos(\pi/x)$. Evaluate $h$ at $x = 1, \tfrac{1}{2},
\tfrac{1}{3}, \tfrac{1}{4}, \dots$ (that is, $x = 1/n$). Show that these values do not all agree,
and use this to argue that $\lim_{x\to 0^+} h(x)$ does not exist. Then explain why \emph{no}
finite table could ever \emph{prove} that a limit \emph{does} exist, even when one does.
\end{exercise}

\begin{exercise}
The Heaviside function $H$ has $\lim_{x\to 0^-}H(x)=0$ and $\lim_{x\to 0^+}H(x)=1$. Define a new
function $S(x) = H(x) + H(-x)$. Work out $S(x)$ explicitly for $x<0$, $x=0$, and $x>0$, then find
both one-sided limits of $S$ at $0$ and decide whether $\lim_{x\to 0}S(x)$ exists. (This is a nice
case where two discontinuous functions combine into something whose two-sided limit \emph{does}
exist.)
\end{exercise}

\begin{exercise}
(\emph{Conceptual / historical.}) In two or three sentences each, describe the distinct
contribution of (a) Newton and Leibniz, (b) Berkeley, (c) Cauchy, and (d) Weierstrass to our
understanding of limits. Your answer to (a) should make clear what they got \emph{right}
intuitively, and your answer to (d) should explain what the shift from ``approaches'' to ``for
every $\eps$ there exists $\delta$'' actually buys us.
\end{exercise}

\begin{exercise}
(\emph{Toward $\eps$--$\delta$.}) For $f(x) = 3x + 2$, we believe $\lim_{x\to 1} f(x) = 5$.
Suppose a skeptic demands that the output be within $0.6$ of $5$, i.e.\ $\abs{f(x) - 5} < 0.6$.
Find an interval of $x$-values around $1$ (of the form $\abs{x - 1} < \delta$ for some $\delta>0$)
that guarantees this. Then redo it for a tolerance of $0.06$. What is the pattern relating the
tolerance to your $\delta$? You have just played the $\eps$--$\delta$ game by hand.
\end{exercise}
